\documentclass[11pt,a4paper]{article}
\usepackage[margin=25mm,headheight=15pt]{geometry}
\usepackage{fontspec}
\setmainfont{texgyrepagella-regular.otf}[BoldFont=texgyrepagella-bold.otf,ItalicFont=texgyrepagella-italic.otf,BoldItalicFont=texgyrepagella-bolditalic.otf]
\setsansfont{texgyreheros-regular.otf}[BoldFont=texgyreheros-bold.otf]
\setmonofont[Path=fonts/,Scale=0.88]{JuliaMono-Regular.ttf}
\usepackage{amsmath,amsthm,amssymb,mathtools}
\usepackage{microtype,xcolor,booktabs,longtable,array,enumitem,fancyhdr,fvextra}
\usepackage[unicode,colorlinks=true,linkcolor=teal!55!black,urlcolor=teal!55!black,citecolor=teal!55!black]{hyperref}
\usepackage{bookmark}
\definecolor{ink}{HTML}{183540}
\definecolor{muted}{HTML}{52626B}
\pagestyle{fancy}\fancyhf{}
\fancyhead[L]{\small\sffamily\color{muted}Uniform Syracuse atom bound}
\fancyhead[R]{\small\sffamily\color{muted}Verified derivation}
\fancyfoot[C]{\small\thepage}
\setlength{\parskip}{5pt}\setlength{\parindent}{0pt}
\setlength{\emergencystretch}{2em}
\newtheorem{theorem}{Theorem}[section]
\newtheorem{lemma}[theorem]{Lemma}
\newtheorem{proposition}[theorem]{Proposition}
\newtheorem{corollary}[theorem]{Corollary}
\theoremstyle{definition}\newtheorem{definition}[theorem]{Definition}
\theoremstyle{remark}\newtheorem{remark}[theorem]{Remark}
\newcommand{\Z}{\mathbb Z}\newcommand{\N}{\mathbb N}\newcommand{\R}{\mathbb R}
\newcommand{\E}{\mathbb E}\newcommand{\PP}{\mathbb P}
\newcommand{\pos}[1]{\left(#1\right)_{+}}
\newcommand{\proj}[2]{\pi_{#1,#2}}
\newcommand{\av}[2]{\langle #2\rangle_{#1}}
\newcommand{\normone}[2]{\lVert #2\rVert_{1,#1}}
\newcommand{\code}[1]{{\small\nolinkurl{#1}}}
\hypersetup{pdftitle={A Uniform Lower Bound for Syracuse Atoms},pdfauthor={Research derivation and Lean verification record}}
\begin{document}
\begin{center}
{\LARGE\bfseries\color{ink}A Uniform Lower Bound\\[4pt]for Syracuse Atoms}\par
\vspace{7pt}
{\large Mathematical derivation and checked Lean source}\par
\vspace{6pt}
{\sffamily Reverification: 12 September 2026}
\end{center}
\begin{abstract}
We derive a constant $c>0$, independent of the depth $n$ and the unit residue
$r\pmod{3^n}$, such that the actual Syracuse probability satisfies
$\mu_n(r)\ge c3^{-n}$. The connection is a positive-part comparison of
filtered backward marks with the full law, followed by a fixed-cylinder
argument and a bounded inverse step. We give the complete derivation of this
connection, state precisely the imported analytic inputs, explain their
summability mechanism, and append the consolidated Lean proof that was
compiled for this document. The argument uses substantial prior work by
Terence Tao and Lech Mazur. It does not prove the Collatz conjecture or
establish historical priority.
\end{abstract}
\tableofcontents
\newpage

\section{The exact theorem and the probability law}
Let $R_n=\Z/3^n\Z$ and $U_n=R_n^\times$. If $h\le q$, write
$\proj hq:R_q\to R_h$ for reduction modulo $3^h$.
Let $A_1,A_2,\ldots$ be independent positive integer variables satisfying
\begin{equation}\label{geom}
 \PP(A_i=v)=2^{-v},\qquad v=1,2,\ldots.
\end{equation}
All powers of $2$ below are inverted in the indicated odd-modulus ring.
Define the probability mass function
\begin{equation}\label{law}
 \mu_n(r)=\PP\left(\sum_{j=1}^{n}3^{j-1}2^{-(A_1+\cdots+A_j)}=r\text{ in }R_n\right).
\end{equation}
The empty sum at level zero is zero. The reversed iid list has the same
law, so this is equivalently the recursive law used in Lean:
\begin{equation}\label{step}
 F_{n,v}(z)=2^{-v}(3\widetilde z+1)\pmod{3^{n+1}},\qquad
 \mu_{n+1}(r)=\sum_{z\in R_n}\sum_{v\ge1}
 \mu_n(z)2^{-v}\mathbf1_{F_{n,v}(z)=r}.
\end{equation}
Here $\widetilde z$ is the canonical integer representative. Changing its
lift by $3^n$ changes $3\widetilde z$ by $3^{n+1}$, so the map is well-defined.
The Lean objects are \code{Erdos1135.Tao.syracStep} and
\code{Erdos1135.Tao.syracPMF}.

\begin{theorem}[Uniform atom lower bound]\label{main}
There exists a real number $c>0$ such that for every integer $n\ge1$ and
every $r\in U_n$,
\begin{equation}\label{mainbound}
 \mu_n(r)\ge \frac{c}{3^n}.
\end{equation}
\end{theorem}
The quantifier order is $\exists c\;\forall n\;\forall r$. No hypothesis
about deterministic Collatz convergence is present.

For $n\ge1$, the law is supported on $U_n$: reducing the sum in
\eqref{law} modulo $3$ leaves $2^{-A_1}\ne0$. There are $2\cdot3^{n-1}$ units.
Consequently, if $m_n=\min_{r\in U_n}\mu_n(r)$, averaging and
Theorem~\ref{main} give
\begin{equation}\label{theta}
 c3^{-n}\le m_n\le\frac32\,3^{-n}.
\end{equation}
In particular $m_n=\Theta(3^{-n})$, and
$\log m_n=-n\log3+O(1)$, hence $m_n=3^{-n+o(n)}$.
This is stronger than the exponential-rate conjecture discussed in
Tao's 2020 note~\cite{tao2020}. The asymptotic interpretation is an elementary
corollary explained here; the literal audited Lean endpoint is
\eqref{mainbound}.

\section{Normalization, projectivity, and the source inputs}
For a function $f:R_q\to\R$, put
\[
 \av qf=3^{-q}\sum_{y\in R_q}f(y),\qquad
 \normone qf=\av q{|f|},\qquad t_+=\max(t,0).
\]
The density used in the imported core library is
\begin{equation}\label{rho}
 \rho_q(y)=\frac23\,3^q\mu_q(y),\qquad \av q{\rho_q}=\frac23.
\end{equation}
This is density relative to unit Haar measure, extended by zero, for
positive $q$. We still take all displayed averages over the \emph{full}
ring $R_q$. At $q=0$, the same definition gives $\rho_0(0)=2/3$.

Reduction in \eqref{law} discards all terms indexed above $h$. Therefore
\begin{equation}\label{projective}
 \sum_{\proj hq(y)=x}\mu_q(y)=\mu_h(x),\qquad
 \sum_{\proj hq(y)=x}\rho_q(y)=3^{q-h}\rho_h(x).
\end{equation}
Each fiber has $3^{q-h}$ points. In particular
$\av q{f\circ\proj hq}=\av hf$.

We use the following imported fine-scale mixing theorem~\cite{tao2019,mazur}:
for every positive integer $A$, there exists $C_A\ge0$ such that
\begin{equation}\label{mix}
 \normone q{\rho_q-\rho_k\circ\proj kq}
 \le \frac{2C_A}{3k^A},\qquad 1\le k\le q.
\end{equation}
It is \code{Tao.taoProp114FineScaleMixing} together with the exact
normalization lemma for the reference density.

The second major input is Mazur's core construction and its uniform
variation estimate. Section~\ref{coreinputs} specifies precisely the facts
used and the quantitative majorant yielding their convergence. These are
\emph{proved library theorems}, not added assumptions in the final Lean
statement. This document does not reproduce the proofs of the entire
imported Fourier and core library; it gives every step from these explicitly
identified results to Theorem~\ref{main}. The supporting source is pinned
and included in the local verification scan.

\section{Prefix transport and the one-sided estimate}
For a finite valuation word $w=(a_1,\ldots,a_d)$, write
$|w|=d$, $s(w)=\sum_i a_i$, and $p(w)=2^{-s(w)}$. Use the outer-first
composition convention
\[
 \Phi_w(z)=\phi_{a_1}\circ\cdots\circ\phi_{a_d}(z),
 \qquad \phi_v(z)=2^{-v}(3z+1).
\]
Thus the constant term of $\Phi_w$ is the finite sum in \eqref{law}.
On $R_{q-d}\to R_q$, $\Phi_w$ is injective: its linear coefficient is
$3^d2^{-s(w)}$. Define the weighted density transport
\begin{equation}\label{transport}
 (T_{w,q}f)(y)=3^d p(w)\sum_{\substack{z\in R_{q-d}\\\Phi_w(z)=y}}f(z).
\end{equation}
Injectivity, followed by summation, gives
\begin{equation}\label{transportnorm}
 \normone q{T_{w,q}f}=p(w)\normone{q-d}f.
\end{equation}
This positive linear map is the mathematical form of the imported
reference prefix transport.

Let $\mathcal W$ be a finite prefix-free family with $|w|\le q$.
Its prefix events are pairwise disjoint under the product law
\eqref{geom}, so
\begin{equation}\label{kraft}
 \sum_{w\in\mathcal W}p(w)\le1.
\end{equation}
Attaching the full residual Syracuse law after each word yields a selected
subdensity of the full law:
\begin{equation}\label{selected}
 S(y):=\sum_{w\in\mathcal W}T_{w,q}\rho_{q-|w|}(y)\le\rho_q(y).
\end{equation}
The factor $2/3$ in \eqref{rho} is the same for all residual lengths,
including zero, so the normalization is consistent.

Suppose $1\le k\le q-|w|$ for all words, and replace each residual density
by its $k$-digit projection. Set
\[
 L(y)=\sum_{w\in\mathcal W}T_{w,q}
       (\rho_k\circ\proj{k}{q-|w|})(y).
\]
By \eqref{transportnorm}, \eqref{mix}, and \eqref{kraft},
\begin{align}
 \normone q{L-S}
 &\le\sum_{w\in\mathcal W}p(w)
    \normone{q-|w|}{\rho_k\circ\proj{k}{q-|w|}-\rho_{q-|w|}}\notag\\
 &\le \frac{2C_A}{3k^A}.
\end{align}
Since $S\le\rho_q$, pointwise $(L-\rho_q)_+\le |L-S|$, whence
\begin{equation}\label{onesided}
 \av q{(L-\rho_q)_+}\le\frac{2C_A}{3k^A}.
\end{equation}
\textbf{The useful change of estimate:} rejected words are favorable to a
lower bound. Their total mass is not an error term in \eqref{onesided}.
This is why the positive part is used instead of absolute error against
the full law.

\section{Filtered kernels and iteration}
For an imported selected prefix family of maximum horizon $H$, a terminal
conductor $k$, and a word filter $0\le\psi(w)\le1$, define
\begin{equation}\label{kernel}
 Kf=\sum_{w\in\mathcal W}\psi(w)T_{w,H+k}
                  (f\circ\proj{k}{H+k-|w|}).
\end{equation}
Positivity gives $(Kf-Kg)_+\le K((f-g)_+)$. The mean identity from
\eqref{transportnorm} gives, for $u\ge0$,
\[
 \av{H+k}{Ku}=\left(\sum_w\psi(w)p(w)\right)\av ku\le\av ku.
\]
Consequently
\begin{equation}\label{contract}
 \av{H+k}{(Kf-Kg)_+}\le\av k{(f-g)_+}.
\end{equation}
Also $K\rho_k\le L$ because $\psi\le1$ and the transported densities
are nonnegative. Absorbing $2/3$ into the constant in \eqref{onesided},
\begin{equation}\label{onestage}
 \av{H+k}{(K\rho_k-\rho_{H+k})_+}\le D_A/k^A.
\end{equation}
The constant is uniform over the root-side parameters and the selected
families. It is not chosen separately at each generation.

For a base $b\ge200$ define $F(b)=b+\lfloor b/100\rfloor$. The imported
horizon is $H(b)=b+\lfloor3b/5\rfloor$. With arbitrary cap and width
schedules, define recursively
\begin{align*}
 q(b,k,0)&=k,&G(b,k,0)&=\rho_k,\\
 q(b,k,N+1)&=H(b)+q(F(b),k,N),&
 G(b,k,N+1)&=K_bG(F(b),k,N).
\end{align*}
In the last line the cap schedule is shifted by one, as in the Lean
recursion. The core word filter is the indicator
\[
 \psi_{b,m}(w)=\mathbf1_{\{b\le |w|+m,\ |w|\le b+m\}}.
\]
The selected family itself is the imported capped least-crossing word
family; it is not selected using the terminal density. Its exact source
name is \code{ndShiftedReferenceSelectedWords}. The results above need its
prefix-free property, bounded horizon, and original probabilities only.

\begin{lemma}[All-generation comparison]\label{iteration}
For every positive integer $A$ there exists $D_A\ge0$ such that
\[
 \av{q(b,k,N)}{(G(b,k,N)-\rho_{q(b,k,N)})_+}
 \le ND_A/k^A\qquad(k\ge1).
\]
\end{lemma}
\begin{proof}
The case $N=0$ has error zero. For the inductive step put
$t=q(F(b),k,N)\ge k$ and split
\[
 (K_bG-\rho_{H(b)+t})_+
 \le (K_bG-K_b\rho_t)_++(K_b\rho_t-\rho_{H(b)+t})_+.
\]
Average and use \eqref{contract}, the induction hypothesis, and
\eqref{onestage}. The result is at most
$ND_A/k^A+D_A/t^A\le(N+1)D_A/k^A$.
\end{proof}

\section{The core inputs and why their tails are uniform}\label{coreinputs}
Fix once and for all
\[
 b_0=32^5,\quad b_{j+1}=F(b_j),\quad
 \operatorname{cap}(j)=\lfloor j/100\rfloor,\quad m(b)=\lceil b^{3/5}\rceil.
\]
For each $N$, let
\[
 k_N=\lfloor b_N/4\rfloor,\quad q_N=q(b_0,k_N,N),\quad
 G_N=G(b_0,k_N,N)
\]
with these cap and width schedules. The use of $k_N$ at every level of
one backward recursion is important: it is a fixed terminal conductor
for that recursion, although it changes when $N$ changes.

Write $s_j$ for the probability that the imported least-crossing word at
floor $b_j$, with cap $\lfloor j/100\rfloor$, lies in the narrow strip
$[b_j-m(b_j),b_j+m(b_j)]$, and set $P_N=\prod_{j<N}s_j$.
The imported core results provide the following facts.
\begin{enumerate}[label=\textbf{C\arabic*.},leftmargin=34pt]
\item $0<s_j\le1$, $\sum_j(1-s_j)<\infty$, and hence
      $P_N\ge p>0$ for every $N$.
\item $G_N\ge0$, it vanishes on nonunits, and
      $\av{q_N}{G_N}=(2/3)P_N$. In particular some unit $y_N$ satisfies
      $G_N(y_N)\ge(2/3)P_N$.
\item For every odd integer $M\ge16^{b_0}$ there exists a real $\ell_M$
      such that
      \begin{equation}\label{uniformtail}
      |\ell_M-G_N(M\bmod3^{q_N})|\le T_N,\qquad T_N\longrightarrow0,
      \end{equation}
      where the \emph{same} sequence $T_N$ works for all such $M$.
\end{enumerate}

Here is the summability mechanism behind C1 and C3, so that the root
uniformity is explicit. Let
\[
 C_{\rm strip}=8\cdot35!\cdot576^{35},\quad g=2013/2000,\quad
 r_1=g(200/201)^6,\quad r_2=g2^{-1/100}.
\]
Both $0<r_1,r_2<1$ are checked in the supporting library. The imported
strip and cap bounds imply
\begin{equation}\label{deficits}
 0\le1-s_j\le \frac{C_{\rm strip}}{b_0^6}(200/201)^{6j}+2^{-j/100}.
\end{equation}
Thus the deficits are summable. Since every $s_j>0$, eventually $s_j\ge1/2$;
then $-\log s_j\le2(1-s_j)$. Therefore
$p=\exp(\sum_j\log s_j)>0$ and $P_N\ge p$. No numerical lower estimate
for $p$ is claimed.

For some fixed $C\ge0$, define
\begin{align}
 B_j={}&\frac23(4b_0^{3/5}+1)6(2^{b_0+1}+16^{b_0})(j+1)^3\notag\\[-2pt]
 &\hspace{8mm}\cdot\left[\frac{2C8^6+C_{\rm strip}}{b_0^6}r_1^j+r_2^j\right].\label{majorant}
\end{align}
Mazur's increment theorem bounds the absolute difference between consecutive
core observables by the original state denominator times $B_j$. For a
singleton root that denominator is exactly $1$, \emph{independently of $M$}.
The singleton observable is exactly $G_j(M\bmod3^{q_j})$, not an approximation
to it. Since $\sum_j(j+1)^3r_i^j<\infty$, we may take
\[
 T_N=\sum_{j\ge N}B_j.
\]
The Cauchy criterion and telescoping prove \eqref{uniformtail}. The
constant $C$, all factors in \eqref{majorant}, and the tail cutoff are
chosen before the physical root. This is the extra information that an
arbitrary $L^1$-mixing family need not possess.

For direct source correspondence, C1 uses
\code{exists_pos_le_coreProbabilityProduct}; C2 uses
\code{exists_unit_rootCoreBackwardMark_ge_product}; C3 uses
\code{exists_coreMarkedLimit_with_root_uniform_tail},
\code{rootCoreSingletonState_denominator_eq_one}, and
\code{rootCoreSingletonState_markedMass_eq_backwardMark}.
These are in the imported core backward-mean and summed-variation modules.

\section{Freezing one positive cylinder}
\begin{lemma}\label{fixed}
There exist $q_*\ge1$, $y_*\in U_{q_*}$, $a>0$, and $N_*$ such that
for every $N\ge N_*$ with $q_N\ge q_*$ and every $z\in R_{q_N}$,
\[
 \proj{q_*}{q_N}(z)=y_*\quad\Longrightarrow\quad G_N(z)\ge a.
\]
\end{lemma}
\begin{proof}
Choose $N_*$ so that $T_j<p/6$ for all $j\ge N_*$. By C2 choose
$y_*\in U_{q_{N_*}}$ with $G_{N_*}(y_*)\ge2p/3$ and put
$q_*=q_{N_*}$, $a=p/3$. For every eligible $M\equiv y_*\pmod{3^{q_*}}$,
\begin{align*}
 G_N(M)&\ge\ell_M-T_N\\
       &\ge G_{N_*}(M)-T_{N_*}-T_N>2p/3-p/6-p/6=p/3.
\end{align*}
All function arguments here are reduced at their respective conductors.
For any fixed finite $z\in R_{q_N}$ above $y_*$, choose an odd natural
representative $M$ of $z$ with $M\ge16^{b_0}$. Such representatives exist
arbitrarily far out because $3^{q_N}$ is odd: among consecutive integer
lifts, parity alternates. Applying the preceding inequality to this $M$
gives the claimed finite-quotient assertion.
\end{proof}
The representative may depend on $N$ and $z$. This causes no dependence
in $a$ or $T_N$, since \eqref{uniformtail} is uniform over every eligible
representative. The favorable residue itself is chosen only once.

\section{Vanishing error and projection onto a fixed subcell}
For $b\ge200$, $\lfloor b/100\rfloor\ge2$. Thus $b_N\ge b_0+2N$.
Using the remainder bound $b_N\le4k_N+3$ gives
\[
 b_0+2N\le4k_N+3,\qquad N\le2k_N,\qquad q_N\ge k_N\ge1.
\]
Lemma~\ref{iteration}, with $A=2$, yields for $N>0$
\begin{equation}\label{error}
 \varepsilon_N:=\av{q_N}{(G_N-\rho_{q_N})_+}
 \le\frac{ND_2}{k_N^2}\le\frac{4D_2}{N}\longrightarrow0.
\end{equation}
Only linear growth is needed here; stronger geometric floor estimates
were used in the imported summability analysis.

\begin{lemma}[Fixed-fiber estimate]\label{fiber}
Let $h\le q$, $x\in R_h$, and suppose $f(z)\ge a$ whenever
$\proj hq(z)=x$. Then
\[
 a\le\rho_h(x)+3^h\av q{(f-\rho_q)_+}.
\]
\end{lemma}
\begin{proof}
For every point of the fiber,
$a\le\rho_q(z)+(f(z)-\rho_q(z))_+$. Sum over that fiber and bound its
nonnegative error by the error over the full ring. By \eqref{projective},
\[
 3^{q-h}a\le3^{q-h}\rho_h(x)+3^q\av q{(f-\rho_q)_+}.
\]
Division by the positive fiber cardinality proves the assertion.
\end{proof}

\begin{proposition}[A positive cylinder for the actual law]\label{actual}
For $q_*,y_*,a$ from Lemma~\ref{fixed}, every $h\ge q_*$ and every
$x\in R_h$ with $\proj{q_*}h(x)=y_*$ satisfy $\rho_h(x)\ge a$.
\end{proposition}
\begin{proof}
Fix $h$ and $x$ \emph{first}. For all sufficiently large $N$, $N\ge N_*$
and $q_N\ge h$, and $G_N\ge a$ on the entire fiber above $x$. Thus
Lemma~\ref{fiber} and \eqref{error} give
\[
 a\le\rho_h(x)+3^h\varepsilon_N\le\rho_h(x)+\frac{4D_2\,3^h}{N}.
\]
Letting $N\to\infty$ proves the result. Equivalently, if
$\delta=a-\rho_h(x)>0$, choose an integer
$N>\max(N_*+2h+1,4D_2 3^h/\delta)$; the same inequality contradicts
$\delta>0$. This finite contradiction is the implementation in Lean.
\end{proof}
The factor is $3^h$, not a harmless uniform constant in $h$. The proof
works because the test conductor $h$ is fixed before $N$ is enlarged.
No assertion is inferred by multiplying a moving-conductor $L^1$ error
by its exponentially growing cardinality.

\section{Bounded inverse coverage of the fixed cylinder}
For $t,c\in\N$ put
\[
 d_t=(4^t-1)/3,\qquad c_t=4^tc+d_t=(3c+1)d_t+c.
\]
These are natural numbers and
\begin{equation}\label{sibling}
 3c_t+1=4^t(3c+1).
\end{equation}
The residues $d_t$, $0\le t<3^q$, permute $R_q$. Indeed, for $s<t$,
\[
 d_t-d_s=4^s\frac{4^{t-s}-1}{3},\qquad
 v_3(d_t-d_s)=v_3(t-s)<q.
\]
The identity $v_3(4^m-1)=1+v_3(m)$ follows by induction on $v_3(m)$:
for $3\nmid u$, the geometric sum $(4^u-1)/3$ is $u\not\equiv0\pmod3$;
and if $x\equiv1\pmod3$, then $x^2+x+1$ has valuation exactly one.
Factor $x^3-1=(x-1)(x^2+x+1)$ repeatedly. This proves distinctness,
and cardinality gives surjectivity. Since $3c+1$ is invertible modulo
$3^q$, the affine formula for $c_t$ also permutes all residues.
The corresponding existing Lean result is \code{CollatzTargetDensity.sibling_residue}.

\begin{lemma}[Bounded inverse valuation]\label{inverse}
If $q\le h$, $y\in R_q$, and $r\in U_{h+1}$, there are
$1\le v\le2\cdot3^q+2$ and $z\in R_h$ with
$\proj qh(z)=y$ and $F_{h,v}(z)=r$.
\end{lemma}
\begin{proof}
Let $R$ be the canonical representative of $r$. If $R\equiv1\pmod3$,
take $e=2$ and $c=(4R-1)/3$; if $R\equiv2\pmod3$, take $e=1$ and
$c=(2R-1)/3$. In both cases $c\in\N$ and $3c+1=2^eR$.
Choose $0\le t<3^q$ so that $c_t\equiv y\pmod{3^q}$.
Then \eqref{sibling} gives $3c_t+1=2^{2t+e}R$. Set $v=2t+e$ and
$z=c_t\pmod{3^h}$. Reducing $c_t$ before multiplying by $3$ preserves
the equation modulo $3^{h+1}$. Division by $2^v$ proves
$F_{h,v}(z)=r$, and $v\le2\cdot3^q+2$.
\end{proof}
No assertion about whether $r$ converges is used. This is finite residue
arithmetic, and the bound depends on the prescribed cylinder only.

\section{Proof of the uniform bound at every depth}
Apply Lemma~\ref{inverse} to the cylinder from Proposition~\ref{actual}.
Set
\[
 V=2\cdot3^{q_*}+2,\qquad d=3\,2^{-V}a>0.
\]
If $h\ge q_*$ and $r\in U_{h+1}$, select the inverse child $z$ above
$y_*$. A single nonnegative term of \eqref{step} yields
\[
 \mu_{h+1}(r)\ge2^{-v}\mu_h(z).
\]
Multiplying by $(2/3)3^{h+1}$, Proposition~\ref{actual} gives
\begin{equation}\label{spread}
 \rho_{h+1}(r)\ge3\,2^{-v}\rho_h(z)\ge3\,2^{-V}a=d.
\end{equation}
This is a lower bound for all units at every depth at least $q_*+1$.
For $1\le h<q_*+1$, fix a unit $r\in R_h$. Every lift of $r$ in
$R_{q_*+1}$ is a unit. By \eqref{projective}, $\rho_h(r)$ is the average
of the densities on that fiber, all at least $d$ by \eqref{spread}.
Therefore $\rho_h(r)\ge d$ at every positive depth.
Finally \eqref{rho} implies
\[
 \mu_h(r)=\frac32\,3^{-h}\rho_h(r)\ge\frac32d\,3^{-h}.
\]
Taking $c=(3/2)d=(9/2)2^{-V}a>0$ proves Theorem~\ref{main}.

\section{Verification, dependency boundaries, and reproduction}
\input{verification-summary.tex}
The consolidated source in Appendix~\ref{appendix} contains all five
modules of the new comparison-to-floor chain in order. Their local import
lines are replaced by three supporting imports at the top. The appended
file was compiled separately from the original modules; it is not a
hand-transcribed shortened proof. Its last command prints the transitive
axioms of the literal theorem.

The library supplies the probability definitions, Fourier mixing, core
construction and variation estimates, as well as sibling arithmetic.
The appendix therefore compiles \emph{inside the pinned project}, not in
an empty Lean installation. It does not reproduce Mathlib or all supporting
source files. The complete original supporting library remains in
\code{ProofAtlasAttack/vendor/positive-density-log-time-collatz}.
The imported source's corrected release commit is
\begin{center}\small\texttt{0aef8b0cfaaf6959500063472f2e328c60df8d54}.\end{center}

From the repository root, reproduce the checks and the document using:
\begin{Verbatim}[fontsize=\small,breaklines=true]
cd ProofAtlasAttack
LEAN_NUM_THREADS=4 python3 scripts/verify.py
lake env lean ../output/pdf/syracuse-uniform-floor/SyracuseUniformFloor.lean
cd ../output/pdf/syracuse-uniform-floor
xelatex -halt-on-error syracuse-uniform-floor.tex
xelatex -halt-on-error syracuse-uniform-floor.tex
\end{Verbatim}
The source is pinned to \code{leanprover/lean4:v4.30.0-rc2}. The package
check uses cached pinned third-party dependencies; it is not an independent
LeanChecker replay or a cold rebuild of all Mathlib. A passing kernel check
establishes the exact formal declaration relative to those foundations;
it does not establish priority or independently referee this exposition.

The bound is existential and gives no practical numerical constant.
It concerns the random variables \eqref{law}. It does not replace the
missing deterministic argument showing that every positive Collatz
trajectory reaches $1$. In particular, density-one convergence, exclusion
of all other cycles, and exclusion of all divergent orbits do not follow
merely by restating this probability estimate.

\begin{thebibliography}{9}
\bibitem{tao2019} T. Tao, \emph{Almost all orbits of the Collatz map attain
almost bounded values}, Forum of Mathematics, Pi 10 (2022), e12.
Preprint arXiv:1909.03562. \url{https://arxiv.org/abs/1909.03562}.
\bibitem{tao2020} T. Tao, \emph{Equidistribution of Syracuse random variables
and density of Collatz preimages}, 25 January 2020.
\url{https://terrytao.wordpress.com/2020/01/25/equidistribution-of-syracuse-random-variables-and-density-of-collatz-preimages/}.
\bibitem{mazur} L. Mazur, \emph{Explicit Positive-Density Collatz Convergence
in Logarithmic Time}, 2026. The local proof uses the pinned, corrected
Apache-2.0 Lean source described above.
\url{https://www.proofatlas.ai/formalizations/positive-density-log-time-collatz/}.
\bibitem{lean} L. de Moura and S. Ullrich, \emph{The Lean 4 theorem prover
and programming language}, CADE 28 (2021), pp. 625--635.
\bibitem{mathlib} The mathlib Community, \emph{The Lean mathematical library},
CPP 2020, pp. 367--381. The exact dependency revisions used here are recorded
in the accompanying verification report.
\end{thebibliography}

\clearpage\appendix
\section{Complete consolidated Lean proof}\label{appendix}
The following listing is the exact UTF-8 content of
\code{SyracuseUniformFloor.lean}. Automatic line wrapping is typesetting
only; use the accompanying source file to run it. The source uses the
three declared supporting imports and includes no unfinished proof steps.
Comments inherited from early stages describe the scope of those individual
modules; the final theorem appears at the end of the complete chain.

